\chapter{Objetivos}
\label{cap:objetivos}

O principal objetivo do projeto diz respeito ao desenvolvimento de um sistema
pericial capaz de utilizar os dados dos sensores simulados de um veículo
autónomo para avaliar situações de risco e apoiar a tomada de decisão mais
adequada em cenários controlados. Neste sentido, foram definidos os seguintes
objetivos específicos:
\begin{itemize}
  \item Identificar os sensores e os dados necessários para apoiar o processo
    de inferência;
  \item Modelar e considerar cenários e elementos rodoviários suficientes para
    dar suporte à tomada de decisão;
  \item Adquirir conhecimento do domínio através da pesquisa e da interação
    com o perito;
  \item Representar o conhecimento através de factos, regras, estados e
    prioridades;
  \item Apresentar uma representação conceptual do funcionamento e do
    processo de decisão do sistema;
  \item Desenvolver, inicialmente, a solução em Drools e preservar a
    equivalência conceptual para a posterior implementação em Prolog;
  \item Permitir a rastreabilidade e explicabilidade das decisões sugeridas e
    descartadas (\emph{How}, \emph{Why}, \emph{Why Not});
  \item Analisar abordagens alternativas ou complementares, como o A* e a
    Teoria dos Jogos;
  \item Explorar a forma como os dados provenientes dos sensores devem ser
    processados e considerados, nos diferentes ciclos de inferência;
  \item Implementar mecanismos de prioridade e de resolução de conflitos
    quando diferentes regras sugerem propostas de ação distintas, em
    simultâneo;
  \item Avaliar situações de risco, através da combinação dos diferentes
    factos e dados disponíveis;
  \item Construir uma base de conhecimento capaz de representar diferentes
    tipos de factos, relações, condições, prioridades e situações relevantes
    para o domínio.
\end{itemize}
